\section{Support shape and weighted variance}\label{sec:weighted-shape}

We now impose the additional third-moment bound
\begin{equation}\label{eq:weighted-resource-range}
                       \beta\le\beta_*:=\frac{18481}{25000}.
\end{equation}
It supplies the strict scalar curvature inequalities proved below.
Throughout these sections we use the representation in
Theorem~\ref{thm:structure}, with data
\eqref{eq:max-data}. Omit deterministic coordinates and put
\[
 H_i=\diam(\supp X_i),\qquad v_i=\E X_i^2,\qquad
 V(h)=\sum_{H_i\le h}v_i,\qquad
 W_s=\sum_i v_iH_i^{-s}.
\]
Constants denoted by $K$ may depend on this fixed representation and
on explicitly fixed exponents, but not on a selected set of its
coordinates or on a scale tending to zero.

\begin{lemma}[The scalar curvature margin]\label{lem:weighted-scalar}
Let $C\ge\ce$, $D=C\beta$, and define
\[
 L_0=\frac{15D+\phi'''(z)}4,\qquad
 c_\varepsilon=3C+\varepsilon\phi''(z)/2,\quad
 \varepsilon\in\{-1,1\}.
\]
Under \eqref{eq:weighted-resource-range},
$c_\varepsilon>0$ and $c_\varepsilon^2>3DL_0$ for both signs.
\end{lemma}
\begin{proof}
Set $r=486833/500000$. The identity
$\phi(0)/\ce=6/(3+\sqrt{10})<r$ follows by squaring the positive
number $6/r-3$. With $x=|z|$, a lower bound for
$(c_\varepsilon^2-3DL_0)/C^2$ is
\[
 \left[3-\frac r2|x^2-1|e^{-x^2/2}\right]^2
 -\frac{3\beta_*r}{4}|3x-x^3|e^{-x^2/2}
 -\frac{45}{4}\beta_*^2.
\]
All square brackets are positive. On $0\le x\le1/2$, replace
$e^{-x^2/2}$ by its upper bound
$T(x)=1-x^2/2+x^4/8-x^6/48+x^8/384$. This gives the polynomial
\[
 P(x)=\left[3-\tfrac r2(1-x^2)T(x)\right]^2
       -\tfrac{3\beta_*r}{4}(3x-x^3)T(x)
       -\tfrac{45}{4}\beta_*^2.
\]
Here is a finite rational positivity check, specified entirely by
these formulas. Write $P(x)=\sum_{j=0}^{20}p_jx^j$ and start with
$b_i=\sum_{j\le i}p_j2^{-j}\binom ij/\binom{20}j$.
Successive midpoint subdivision replaces a coefficient row by the
left and right edges of its averaging triangle. On the five indicated
intervals the smallest resulting coefficient exceeds the stated
rational number:
\[
\begin{array}{c|ccccc}
 I &[0,1/8]&[1/8,3/16]&[3/16,7/32]&[7/32,1/4]&[1/4,1/2]\\ \hline
 \min b_i(I)>&1/40&1/1000&1/30000&1/1000&9/1000 .
\end{array}
\]
The corresponding subdivision words are $00,010,0110,0111,1$.
The Bernstein basis is nonnegative and sums to one, so these finite
rational comparisons prove $P>0$ on $[0,1/2]$.
For $x\ge1/2$, elementary differentiation gives
$|\phi''(x)|/\phi(0)<2/3$ and
$|\phi'''(x)|/\phi(0)<7/5$. The remaining lower bound is
\[
 (3-r/3)^2-\tfrac{3\beta_*r}{4}\tfrac75
                  -\tfrac{45}{4}\beta_*^2
       =\frac{35773617637}{140625000000}>0.
\]
Finally $3C-\phi(0)/2>0$ proves positivity of both $c_\varepsilon$.
\end{proof}
\begin{lemma}[Distances of the extreme support points from zero]\label{lem:weighted-shape}
There is $\kappa>0$ such that every sufficiently late genuine
three-point coordinate, with extreme states $-a_i,b_i$, satisfies
\[
 \min(a_i,b_i)\ge\kappa H_i.
\]
Consequently, after enlarging a fixed constant to include the finite
prefix,
\begin{equation}\label{eq:weighted-shape-bound}
                         H_i\E|X_i|\le K v_i .
\end{equation}
For a two-point coordinate the latter identity holds with $K=2$.
\end{lemma}
\begin{proof}
For a genuine triple $x<y<t$ in $[-h,h]$, vary its raw mean and
second moment, keeping its support fixed. The probability map is
nonsingular because its Vandermonde determinant is nonzero. If $r$
is its centered third absolute moment and
$k=[u|u|]_{x,y,t}$, divided differences give
\[
 |k|\le1,\quad |r_\mu|\le12h^2,\quad |r_v|\le3h,\quad
 |r_{\mu\mu}|\le30h,\quad r_{\mu v}=-3k,\quad r_{vv}=0.
\]
For example the last two identities follow by differentiating
$\E|Y-\mu|^3$ and using the mass, first-moment and second-moment
coordinates; the second divided difference of $u|u|$ lies between
$-1$ and $1$. The exact objective second variation, obtained as in
\eqref{eq:objective}, is
\[
 (-3D-Cr_{\mu\mu})e^2+
   (\phi''(z)+6Ck+3Cr_\mu)ef+
   (-L_0+3Cr_v)f^2.
\]
It is nonpositive for all real $e,f$. Since $h\to0$, any limit of
$k$ satisfies $(3Ck+\phi''(z)/2)^2\le3DL_0$. By
Lemma~\ref{lem:weighted-scalar}, $k$ stays a fixed distance from both
$-1$ and $1$ on late triples.

Write their extremes as $-a,b$. For $y\ge0$ and $y\le0$ respectively,
\[
 k=1-\frac{2a^2}{(a+y)(a+b)},\qquad
 k=-1+\frac{2b^2}{(a+b)(b-y)}.
\]
These formulas and centering give the asserted bound on both extremes.
If the middle point is nonnegative, the negative extreme alone
carries $\E|X|/2$ in first moment; hence
$v\ge a\E|X|/2\ge\kappa H\E|X|/2$. Use the positive extreme when
the middle point is nonpositive. For $X=H(B_p-p)$,
$\E|X|=2p(1-p)H$ and $v=p(1-p)H^2$.
\end{proof}

\section{Third-moment growth under grouping and normal comparison}
\label{sec:weighted-grouping}

A group is a sum of disjoint original coordinates. Its cost is the
increase in the sum of their third absolute moments. This cost is
essential even though grouping does not change the full sum law.

\begin{lemma}[Two grouping estimates]\label{lem:weighted-grouping}
For a centered independent family with total variance $V$,
\begin{equation}\label{eq:weighted-variance-cost}
 0\le\E|\textstyle\sum X_i|^3-\sum\E|X_i|^3\le2V^{3/2}.
\end{equation}
For independent raw variables $Y_i$, let
$X_i=Y_i-\E Y_i$, $M=\sum\E|Y_i|$, and $A_-=\sum\E(Y_i)_-$.
Then their grouping cost is at most $2M^3+6A_-V$.
If the coordinates are divided into $B$ classes, with centered
first-moment and variance sums $M_j,V_j$, there is, for every integer
$J\ge1$, a partition into at most $8J+6B$ groups whose cost is at most
\begin{equation}\label{eq:weighted-annular-cost}
 \frac{18M^3}{J^2}+\frac{9S^2}{J},
 \qquad M=\sum_jM_j,\quad S=\sum_j\sqrt{M_jV_j}.
\end{equation}
All bounds pass to countable groups with finite summed variance and
third-moment sum.
\end{lemma}
The first upper bound is a classical third-moment estimate.
Pinelis~\cite[equation (5)]{Pinelis2013} gives the sharper coefficient
$\E|Z|^3=2\sqrt{2/\pi}$ in place of $2$. The elementary proof below
suffices for all the rates used here.
\begin{proof}
The function $\E|x+Y|^3-|x|^3$ is convex for centered $Y$,
because its second derivative is $6(\E|x+Y|-|x|)\ge0$.
Jensen proves the lower bound. The pointwise Taylor bound
\[
 |x+y|^3\le|x|^3+3x|x|y+3|x|y^2+|y|^3
\]
gives an upper increment $3v_i\sqrt{\sum_{j<i}v_j}$ at step $i$.
Summing and comparing with $\int_0^V3\sqrt u\,du$ proves
\eqref{eq:weighted-variance-cost}. Applied to tail differences, the
same inequality proves convergence in $L^3$.

For centered $X$ define
$k_X(t)=\E(X-t)_+-(-t)_+$. It is nonnegative, has supremum
$\E|X|/2$, and integral $v_X/2$. Taylor's integral remainder yields
the exact identity
\[
 \E|X+Y|^3-\E|X|^3-\E|Y|^3
                  =12\int_{\R}k_X(t)k_Y(-t)\,dt.
\]
For a raw variable with first-moment data $M_i,A_i$,
$k_{X_i}(t)\le M_i$ and
$k_{X_i}(-t)\le(M_i-t)_++A_i$ for $t\ge0$. Split the integral at
zero and apply the same bounds to each preceding partial sum. The
increment is at most
\[
 6(M_{<i}M_i^2+M_iM_{<i}^2)+6(A_iV_{<i}+A_{<i}v_i).
\]
The first terms telescope to at most $2M^3$ and the others sum to
at most $6A_-V$.

More generally choose raw laws and signs with
$X_i=\varepsilon_i(Y_i-\E Y_i)$ and keep the two signs in separate
groups. Retain individually every coordinate with
$M_i>M/J$ or $A_i>A_-/J$. Greedily accumulate the remaining weights
$M_i/M+A_i/A_-$ until they reach $1/J$, omitting the second weight
when $A_-=0$. Every completed group has both normalized weights
at most $3/J$. There are at most $4J+2$ groups, including the
retained coordinates and the final group in each sign. The total
cost is at most $18M^3/J^2+18A_-V/J$.

In class $j$ take $Y_i=X_i$, so $A_-=M_j/2$, and allocate
$J_j=\lceil J(M_j/M+\sqrt{M_jV_j}/S)\rceil$.
Then $\sum J_j\le2J+B$ and the two terms sum to
\eqref{eq:weighted-annular-cost}. Upper budgets in place of $M,S$
only decrease the sum of the allocated weights. These observations
also justify removal of any selected original coordinates before
grouping.
\end{proof}

\begin{lemma}[A quadratic majorant for a sum of coordinates]
\label{lem:weighted-majorant}
Let $U$ be any finite or countable sum of original coordinates,
independent of a remainder with CDF $G$, and let $\E U^2=\sigma^2>0$.
There is a quadratic polynomial $P(u)=a+bu+cu^2$ such that
\begin{equation}\label{eq:weighted-majorant}
 Q_U(u):=P(u)+C|u|^3-G(z-u)\ge0,\qquad
                        \E Q_U(U)\le2C\sigma^3.
\end{equation}
If $\sigma\to0$, the original coordinate radii are at most $H$, and
$H/\sigma\to0$, its coefficients can be chosen with
\begin{equation}\label{eq:weighted-coefficients}
 |a|+|b|+|c|\le K,\qquad b+\phi(z)=O(\sigma),\qquad c+A=O(\sigma).
\end{equation}
In particular $Q_U(u)\le K(1+|u|^3)$ globally.
\end{lemma}
\begin{proof}
Put $f(u)=G(z-u)-C|u|^3$ and maximize $\E f(Y)$ at prescribed
first and second moments $(m,s)$ with $s>m^2$. The value is a
finite concave function: $f\le1$, and a two-point law gives a finite
lower bound. A supporting plane at $(0,\sigma^2)$ gives a quadratic
majorant of $f$ at every real point. To include point masses without
a continuity assumption, mix $\delta_x$ with a law approaching the
value at $(0,\sigma^2)$, apply the plane inequality at the interior
moment pair, and cancel its mixing weight. This yields $P(x)\ge f(x)$.
Replacing the original block by any such $Y$ is an admissible
one-coordinate comparison. If its third-moment sum is $\beta_I$,
maximality bounds the value by $F(z)-C\beta_I$, whereas the actual
grouped law has value $F(z)-C\E|U|^3$. Their difference is at most
$C(\E|U|^3-\beta_I)\le2C\sigma^3$, proving
\eqref{eq:weighted-majorant}.

We give the coefficient argument because it fixes the quadratic term
used later. Deleting the original block and applying the universal
inequality to its remainder gives, for $\sigma^2\le1/2$,
\[
 G(z-u)\le\Phi((z-u)/\sqrt{1-\sigma^2})
                    +C(\beta-\beta_I)(1-\sigma^2)^{-3/2}.
\]
Normal Taylor bounds imply
$G(z-u)-F(z)+\phi(z)u\le K(u^2+\sigma^2)$.
Its expectation is zero, so
\begin{equation}\label{eq:weighted-affine-deletion}
             \E|G(z-U)-F(z)+\phi(z)U|\le K\sigma^2.
\end{equation}
For any bounded score $\psi$ with $|\psi|\le1$,
$\E\psi(U)=\E U\psi(U)=0$, replace the whole block by
$(1+t\psi)\Law(U)$, $|t|\le1$. Write
$g=\E G(z-U)\psi$, $d=\E U^2\psi$, $e=\E|U|^3\psi$.
Its total variance and CDF are $1+td,F(z)+tg$, and its third-moment
sum is $\beta+w+te$, where $0\le w\le2\sigma^3$.
Here $|g|\le K\sigma^2$, $|d|\le\sigma^2$. Taylor expansion of
$B^{3/2}[f-\Phi(z/\sqrt B)]$ at $(1,F(z))$, whose $B$ derivative
is $A$, has remainder $O(\sigma^4)$ for $|t|\le1$.
Using both signs of $t$ gives
\[
 |\E\{G(z-U)+AU^2-C|U|^3\}\psi|
                                      \le Cw+K\sigma^4.
\]
The $L^1$ quotient duality for the subspace spanned by $1,U$
therefore supplies $a_0,b_0$ with
\[
 \E|G(z-U)-C|U|^3-a_0-b_0U+AU^2|=O(\sigma^3).
\]
We may take $a_0=F(z)+A\sigma^2-C\E|U|^3$ by changing the error
constant. Equation \eqref{eq:weighted-affine-deletion} and
$\E|U|\ge c\sigma$ give $b_0+\phi(z)=O(\sigma)$.
The latter lower bound follows from convergence of $U/\sigma$ to
normal; the classical Berry--Esseen bound is at most $K H/\sigma$.

Subtract this polynomial from the majorant in
\eqref{eq:weighted-majorant}. Its $L^1$ norm under $U$ is
$O(\sigma^3)$. On polynomials of degree at most two, the normal
$L^1$ norm restricted to $[-2,2]$ controls all three coefficients.
The same bound holds uniformly under $U/\sigma$ once its
Kolmogorov distance to normal is small: absolute quadratics on the
unit coefficient sphere have uniformly bounded variation there.
It follows that
$a-a_0=O(\sigma^3)$, $b-b_0=O(\sigma^2)$, $c+A=O(\sigma)$,
as required.
\end{proof}

\begin{lemma}[Weighted normal comparison]\label{lem:weighted-normal}
In Lemma~\ref{lem:weighted-majorant}, suppose
$(H/\sigma)\log^{3/2}(1/\sigma)\to0$. Then
\begin{equation}\label{eq:weighted-normal-slack}
                         \E Q_U(\sigma Z)\le K(\sigma^3+H),
                         \qquad Z\sim N(0,1).
\end{equation}
If $H/\sigma=O(\sigma^\zeta)$ for a fixed $\zeta>1$, then for every
$2<q<\min(\zeta+1,3)$ and fixed $T$,
\begin{equation}\label{eq:weighted-normal-window}
 \sup_{|t|\le T\sigma}
 |F(z+t)-F(z)-\phi(z)t+At^2|=O_{q,T}(\sigma^q).
\end{equation}
\end{lemma}
\begin{proof}
We record the weighted approximation estimate used in the proof.
For a variance-one sum $W$ of centered independent variables bounded
by $\epsilon$, and $\epsilon(1+L^3)$ sufficiently small,
\begin{equation}\label{eq:weighted-MD}
 |\Pp(W\le x)-\Phi(x)|+|\Pp(W<x)-\Phi(x)|
       \le K\epsilon(1+x^2)\phi(x),\qquad |x|\le L.
\end{equation}
Here is the bounded-array exponential-tilting proof. If
$K(t)=\log\E e^{tW}$, centered summands under the tilted product
law are bounded by $2\epsilon$, so
$|K'''(t)|\le2\epsilon K''(t)$ and
$e^{-2\epsilon|t|}\le K''(t)\le e^{2\epsilon|t|}$.
For $x\ge0$ choose $K'(t)=x$. Integration gives
$t=x+O(\epsilon x^2)$, $s^2=K''(t)=1+O(\epsilon x)$, and
$K(t)-tx=-x^2/2+O(\epsilon x^3)$. Berry--Esseen under the
tilted law, applied to $e^{-tsu}\ind_{\{u\ge0\}}$ or its strict
version, gives
\[
 \Pp(W\ge x)=e^{K(t)-tx}\{J(ts)+O(\epsilon)\},\quad
 J(a)=\int_0^\infty e^{-au}\phi(u)\,du.
\]
The same formula holds with $>$. Since
$J(a)\ge c/(1+a)$, $|(\log J)'(a)|\le\sqrt{2/\pi}$, and
$e^{-x^2/2}J(x)=1-\Phi(x)$, the relative tail error is
$O(\epsilon(1+x^3))$. Mills' bound and reflection prove
\eqref{eq:weighted-MD}. Finite truncations and thresholds tending
to $x$ from both sides prove the countable version. This is the
classical Cram\'er tilting argument; no nonlattice assumption enters.
For a modern independent-array formulation and its conjugate-law
proof, see Beknazaryan, Sang and Xiao~\cite[Theorem 2.1]{BeknazaryanSangXiao2019}.
Here the bounded-summand estimate above supplies the hypotheses directly.

Put $Q_\sigma(x)=Q_U(\sigma x)$, $w(x)=(1+x^2)\phi(x)$, and
$L=\sqrt{20\log(1/\sigma)}$. In the sense of measures on
$[-L,L]$,
\[
 dQ_\sigma=a_\sigma'(x)\,dx+d\mu(x),\quad d\mu\ge0,\quad
 a_\sigma(x)=P(\sigma x)+C\sigma^3|x|^3.
\]
The positive measure comes from $-G(z-\sigma x)$. Bounded
coefficients give $\int w|a_\sigma'|\le K\sigma$.
Since $Q_\sigma\ge0$, integration by parts against $w$ yields
\[
 \int_{-L}^L w\,|dQ_\sigma|
 \le K\{\sigma+(1+L^3)I+(1+L^2)\phi(L)\},\quad
 I=\int_{-L}^L Q_\sigma(x)\phi(x)\,dx.
\]
Both endpoint terms are included. Stieltjes integration by parts
and \eqref{eq:weighted-MD}, with $\epsilon=H/\sigma$, now give
\[
 I\le\E Q_U(U)+K\epsilon
          \{\sigma+(1+L^3)I+(1+L^2)\phi(L)\}.
\]
The term containing $I$ is absorbed. The global cubic bound on
$Q_U$ pays the normal tails and proves
\eqref{eq:weighted-normal-slack}. The left-continuous function
$G(z-\sigma x)$ and its positive jump measure can be used throughout;
the two endpoint versions in \eqref{eq:weighted-MD} justify the
same argument when either law has atoms.

For a fixed $|t|\le T$, the normal density ratio on $[-L,L]$
under translation by $t$ is at most $e^{TL+T^2/2}$.
The preceding weighted-variation argument applied to
$Q_\sigma(x-t)$, together with
\[
 \Pp(|U|>\sigma L)
 \le2\exp\{-L^2/[2(1+\epsilon L/3)]\}
                       =\sigma^{10-o(1)},
\]
therefore gives
\begin{equation}\label{eq:weighted-shifted-slack}
 \sup_{|t|\le T}\E Q_U(U-\sigma t)
                         \le\sigma^{-o(1)}(\sigma^3+H).
\end{equation}
The omitted actual tail is bounded by Cauchy--Schwarz and
$\E|U|^6=O(\sigma^6)$; all fixed normalized moments are bounded
by the elementary bounded-summand exponential-moment estimate.

The coordinate contact identities determine the linear coefficient more
accurately. Multiplying each identity by $X_i$ and summing gives
\[
 \E[UG(z-U)]=-\phi(z)\sigma^2-A\E U^3+C\sum_i\eta_i,\quad
 \eta_i=\E[X_i|X_i|^3]-3v_i\E[X_i|X_i|].
\]
Consequently
\[
 (b+\phi(z))\sigma^2
 =-(c+A)\E U^3-C\E[U|U|^3]
                 +\E[UQ_U(U)]+C\sum_i\eta_i.
\]
Here $|\E U^3|\le H\sigma^2$, $\sum|\eta_i|\le4H^2\sigma^2$.
Splitting at $|U|=\sigma\sqrt{M\log(1/\sigma)}$ and using
nonnegativity, $\E Q_U(U)=O(\sigma^3)$ and bounded-summand tails
gives $\E|U|Q_U(U)=O(\sigma^4\sqrt{\log(1/\sigma)})$.
Thus $b+\phi(z)=O(\sigma^2\sqrt{\log(1/\sigma)})$ under the
hypothesis $\zeta>1$. Finally the exact identity
\[
 F(z+\sigma t)-F(z)
 =-b\sigma t+c\sigma^2t^2
 +C\{\E|U-\sigma t|^3-\E|U|^3\}
 +\E Q_U(U)-\E Q_U(U-\sigma t)
\]
and \eqref{eq:weighted-shifted-slack} prove
\eqref{eq:weighted-normal-window}.
\end{proof}

\subsection{The closed-CDF phase identity}
\label{subsec:weighted-phase}
We use two elementary exponential bounds for independent centered
variables. If their ranges have lengths $d_i$, integration of the
tilted log-MGF bound $K_i''(t)\le d_i^2/4$ gives
$\E e^{t\sum Y_i}\le\exp(t^2\sum d_i^2/8)$, the bounded-range
estimate of Hoeffding~\cite{Hoeffding1963}. If instead $|Y_i|\le H$,
the exponential series and $k!\ge2\cdot3^{k-2}$ for $k\ge2$ give
\[
 \log\E e^{t\sum Y_i}
 \le\frac{t^2\sum\E Y_i^2}{2(1-H|t|/3)},\qquad H|t|<3.
\]
Chernoff's inequality gives the Bernstein bound used above and below.
It also gives bounded moments of every fixed order after normalization
when $H/\sqrt{\sum\E Y_i^2}$ is bounded.

Unbiased rounding to the adjacent points of a lattice of mesh $\ell$
preserves the mean and increases variance by at most $\ell^2/4$.
For $N$ independent coordinates, choose the first $N-1$ lattice phases
uniformly and the last so that their sum is $z$ modulo $\ell$.
For each fixed phase vector the rounded coordinates are independent.
If their sum is $L_{\boldsymbol u}$ and their unrounded sum has CDF $F$,
then
\begin{equation}\label{eq:weighted-phase-identity}
 \E_{\boldsymbol u}F_{L_{\boldsymbol u}}(z)
       =\E F(z-W_\ell),\quad
 W_\ell=T_1+\cdots+T_{N-1}-U_\ell,
\end{equation}
where the $T_i$ are independent with density
$\ell^{-1}(1-|x|/\ell)_+$ and $U_\ell$ is uniform on $[0,\ell]$.
To verify the sign and endpoints, the phase-averaged mass at
$z+k\ell$ is $\ell$ times the convolution of the $N$ triangular
kernels with the original law. For the last kernel,
$\ell\sum_{k\le0}k_\ell(z+k\ell-y)$ equals one for $y\le z$,
$(z+\ell-y)/\ell$ for $z<y<z+\ell$, and zero for $y\ge z+\ell$.
Tonelli proves \eqref{eq:weighted-phase-identity}, including atoms.
In particular
\[
 \E W_\ell=-\ell/2,\qquad \E W_\ell^2=(N+1)\ell^2/6.
\]
For $N\ge3$ every pair of phases is independent. If $v_{\boldsymbol u}$
is their total rounding variance, then
\begin{equation}\label{eq:weighted-phase-variance}
 \E v_{\boldsymbol u}=N\ell^2/6,\quad
 \Var(v_{\boldsymbol u})\le N\ell^4/64,\quad
 |\operatorname{Cov}(v_{\boldsymbol u},F_{L_{\boldsymbol u}}(z))|
                                                    \le\sqrt N\ell^2/16.
\end{equation}
These follow by integrating $\ell^2t(1-t)$ over a uniform fractional
part and using $0\le F_{L_{\boldsymbol u}}\le1$.

\begin{lemma}[Grouping at a selected normal window]
\label{lem:weighted-window-phase}
Suppose \eqref{eq:weighted-normal-window} holds with $2<q<3$ at
scales $\sigma\downarrow0$. Suppose the original sum is partitioned
into $N\to\infty$ independent groups with actual cost $w\ge0$.
The maximum is impossible if, with $L=\log(1/\sigma)$,
\begin{equation}\label{eq:weighted-window-costs}
 N\sigma^{2q-2}L\to0,\qquad
       w\sqrt{NL}/\sigma\to0,\qquad wN^{2/3}\to0.
\end{equation}
\end{lemma}
\begin{proof}
The conditions permit a mesh
\[
 \max(w,\sigma^q)\ll\ell\ll
                   \min(\sigma/\sqrt{NL},N^{-2/3}).
\]
Indeed $\sigma^qN^{2/3}\to0$ follows from the first condition and
$q<4$. The triangular noise in
\eqref{eq:weighted-phase-identity} leaves $[-\sigma,\sigma]$ with
probability at most $\exp\{-c\sigma^2/(N\ell^2)\}$; its first two
tail moments satisfy the same bound up to polynomial factors.
This is $o(\ell)$ because $\ell\ge\sigma^q$ and the exponent
dominates $L$. Thus the window gives an average CDF increment
$\phi(z)\ell/2-A(N+1)\ell^2/6+O(\sigma^q)+o(\ell)$.

For a fixed phase, the normalized objective expands as
\[
 (1+v)^{3/2}[F_{\boldsymbol u}-\Phi(z/\sqrt{1+v})]
 =F_{\boldsymbol u}-\Phi(z)+Av
       +\tfrac32(F_{\boldsymbol u}-F(z))v+O(v^2).
\]
Equations \eqref{eq:weighted-phase-variance} cancel the term
$-AN\ell^2/6$. The remaining covariance and normalization errors are
$O(\ell^2+\sqrt N\ell^2+N\ell^3+N^2\ell^4)+o(\ell)$.
The actual cubic interpolation cost is at most
$K\ell^2\sum_i\E|Z_i|+KN\ell^3
\le K\sqrt N\ell^2+KN\ell^3$, since the groups have total
variance one. Every error and $w$ is $o(\ell)$. The average
objective is therefore positive, a contradiction.
\end{proof}

\begin{lemma}[Small-diameter Fourier alternative]\label{lem:weighted-Fourier}
If $V(r)\ge c r^p$ for all sufficiently small $r$, for some fixed
$p<2$, the original representation cannot maximize.
\end{lemma}
\begin{proof}
We give the approximation needed for this implication; smoothness of
the original density alone would not suffice. Put
$w(t)=\exp\{-t^2V(1/(8t))/12\}$ for $t\ge1$. The hypothesis
makes $\int_1^\infty t^j w(t)\,dt$ finite for every fixed $j$.
For a centered coordinate of diameter $d\le1/(8t)$ and variance $v$,
each of the following characteristic functions has modulus at most
$e^{-vt^2/12}$, for $1\le t\le\pi/\ell$:
the original law, its accompanying compensated Poisson law,
its unbiased rounding to $\ell\Z$, and the following rounding of its
Poisson law. Interpolate the measure $x^2\Law(X)(dx)$ to $\ell\Z$;
interpret its mass at zero as Gaussian variance and its mass at
$u\ne0$ as $u^2$ times a Poisson intensity.

For the original and Poisson versions use
$1-\cos s\ge s^2/4$ on $|s|\le1$, with an independent copy for the
original law. If $t\ell\le1/4$, the directly rounded diameter is at
most $d+2\ell$, so the same estimate applies. Otherwise $d<\ell/2$,
and the directly rounded law has masses $\rho/2,\rho/2$ at
$-\ell,\ell$, with $\rho=\E|X|/\ell<1/2$.
Since $\E|X|\ge v/d\ge8tv$ and
$1-\cos s\ge2s^2/\pi^2$ on $[0,\pi]$, its damping is stronger than
the stated bound. For the interpolated Poisson measure,
$|tu|\le\pi+1/8$ and $1-\cos(tu)\ge t^2u^2/12$; the Gaussian
mass at zero also has the required damping. Omitting one contributing
factor loses at most $e^{1/768}$.

Order the original variances decreasingly, put
$A_k=\sum_{i>k}v_i^2$, and choose $\ell_k=\sqrt{A_k/k}$.
Then $A_k=o(\ell_k)$ and $k\ell_k\to0$. Replace the tail by its
accompanying Poisson law, interpolate its second-moment measure as
above, and round the first $k$ coordinates directly. The three full
characteristic-function errors, with damping from all remaining
factors, are bounded respectively by
\[
 K A_k t^4w(t),\qquad K\ell_k^2t^4w(t),\qquad
                         Kk\ell_k^2t^2w(t).
\]
For the first use $|e^{f_i-1}-f_i|\le v_i^2t^4/8$.
For the second use
$(e^{itx}-1-itx)/x^2=-t^2\int_0^1(1-s)e^{istx}\,ds$,
whose second derivative is bounded by $t^4/12$.
The interpolation preserves both total variance and the Poisson contribution
to the additive third-moment quantity, because $|x|$ is affine between adjacent lattice
nodes, including those adjoining zero. Direct rounding costs
$O(k\ell_k^2)$ in variance and third-moment sum by bounded original
supports.

The created Gaussian variance $a_k$ is at most $\sum_{i>k}v_i\to0$.
Replace it by $n$ fair centered Bernoullis of span $\ell_k$, taking
$n$ nearest to $4a_k/\ell_k^2$. Their variance differs by
$O(\ell_k^2)$ and their total third-moment sum is $O(a_k\ell_k)=o(\ell_k)$.
For $|t\ell_k|\le\pi$,
$|\cos(t\ell_k/2)|^n\le e^{\pi^2/16}e^{-a_kt^2/2}$;
on $|t\ell_k|\le1$, logarithmic expansion gives an additional full
error $K\ell_k^2(t^2+t^4)w(t)$. Outside that range use the damping
itself. The final admissible law $Z_k$ lies on a translate of
$\ell_k\Z$, has variance $1+o(\ell_k)$ and additive third-moment quantity
$\beta+o(\ell_k)$, and satisfies
\[
 \int_0^{\pi/\ell_k}
     |\widehat Z_k(t)-\widehat S(t)|\,\frac{dt}{t}=o(\ell_k),
 \qquad
 \int_0^{\pi/\ell_k}|\widehat Z_k(t)-\widehat S(t)|\,dt=o(1).
\]
The bounds below $t=1$ use $w=1$. All passages through infinite
products are justified by the displayed summed variance bounds.

Fourier inversion gives a smooth density $f$ for $S$.
At a lattice point $x$, the exact mid-CDF formula is
\[
 F_{Z_k}(x)-\tfrac12\Pp(Z_k=x)
 =\tfrac12-\frac{\ell_k}{2\pi}\int_0^{\pi/\ell_k}
       \Im(e^{-itx}\widehat Z_k(t))\cot(t\ell_k/2)\,dt.
\]
The inequalities
$0\le(\ell/2)\cot(t\ell/2)\le1/t$ and
$|1/t-(\ell/2)\cot(t\ell/2)|\le K\ell^2t$ yield
$\Pp(Z_k=x)=\ell_k f(x)+o(\ell_k)$ and
$F_{Z_k}(x)=F(x)+\ell_k f(x)/2+o(\ell_k)$ uniformly at lattice
points. At the maximizing threshold $f(z)=\phi(z)>0$.
Choose a lattice point within $\ell_k/2$ of $z$; its original
discrepancy differs from $D$ by $O(\ell_k^2)$. After the variance
and third-moment corrections, its objective has gain
$\phi(z)\ell_k/2+o(\ell_k)>0$.
\end{proof}

\section{The first weighted variance moment}\label{sec:weighted-first}

\begin{lemma}\label{lem:weighted-subunit}
Under \eqref{eq:weighted-shape-bound}, $W_{1/2}$ and $W_{3/4}$ are finite.
\end{lemma}
\begin{proof}
For a cutoff $t$ set $M(t)=\sum_{H_i>t}\E|X_i|$ and
\[
 {\cal K}(t)=\sum_i v_i\min(1,t/H_i)
              =V(t)+t\sum_{H_i>t}v_i/H_i.
\]
Then $M(t)\le K{\cal K}(t)/t$ and $V(t)\le{\cal K}(t)$.
The exact Tonelli identity is
\begin{equation}\label{eq:weighted-Tonelli}
 \sum_{H_i\le1}v_iH_i^{-s}
                 =V(1)+s\int_0^1 V(u)u^{-s-1}\,du.
\end{equation}

We first prove $W_{1/2}<\infty$. If it diverges, good points with
$V(u)\ge u^a$, $a=13/24$, occur arbitrarily near zero.
By Lemma~\ref{lem:weighted-Fourier} there are also
$r\downarrow0$ with $V(r)\le r^p$, $p=15/8$.
Take the largest good point below $r$, as a supremum. Monotonicity
and right-continuity of the closed variance function give exactly
\[
 V(h)=h^a,\quad V(u)<u^a\ (h<u<r),\quad h\le r^{p/a}.
\]
There is a gap below $r$ because $V(r)<r^a$; approaching $h$
from above and below proves equality even at a variance atom.
Set $q=23/8$, $t=h^{37/72}$. Then $h<t<r$, and integration
by parts for the actual discrete variance measure gives
\[
 {\cal K}(t)\le t/r+t^a/(1-a).
\]
Indeed the part with $H_i\le r$ equals
$tV(r)/r+t\int_t^r V(u)u^{-2}\,du$.
Select finitely many coordinates with $H_i\le h$ until their
variance first reaches $h^{2/3}$. This is possible since
$h^a>2h^{2/3}$; the overshoot is at most $h^2$. Its actual
standard deviation is $\sigma\sim h^{1/3}$ and
Lemma~\ref{lem:weighted-normal} gives the window with exponent $q$.
The two positive margins are
\[
 \sigma^{-2/3}{\cal K}(t)
       \le K(h^{1/360}+h^{97/1728}).
\]
Thus, with $L=\log(1/\sigma)$,
$(1+M(t))\sigma^{q-2}L\to0$ and $V(t)^3L/\sigma^2\to0$.

To check the grouping consequence explicitly, take
$J=\lceil(1+M(t))\sigma^{-q}\rceil$ and retain the first
$\lfloor\sqrt J\rfloor$ original coordinates separately.
Group the remaining large coordinates by
Lemma~\ref{lem:weighted-grouping}, and the small ones by variance:
retain coordinates exceeding $V(t)/J$ and accumulate the others
into groups of variance at most $2V(t)/J$. The complete count is
$N=O(J)$, $N\to\infty$, and
\[
 w\le K\{(1+M)^3/J^2+(1+M)/J+V^{3/2}/\sqrt J\}.
\]
The three expressions in \eqref{eq:weighted-window-costs} tend to
zero. For the second, the middle term is bounded by
$K\sqrt{(1+M)\sigma^{q-2}L}$, the first is smaller by
$(1+M)\sigma^q$, and the last is $KV^{3/2}\sqrt L/\sigma$.
For the third, the middle term is
$((1+M)\sigma^{q-2}L)^{2/3}\sigma^{(4-q)/3}L^{-2/3}$,
and the same two comparisons apply. This contradicts
Lemma~\ref{lem:weighted-window-phase}.

Now write $W=W_{1/2}<\infty$ and suppose $W_{3/4}=\infty$.
Repeat the crossing argument with
\[
 a=7/8,\quad p=15/8,\quad Q_0=15/7,\quad q=29/14,\quad t=\sqrt h.
\]
It gives $V(h)=h^a$, $V(u)<u^a$ on $(h,r)$ and
$h\le r^{p/a}$. Select original variance
$\sigma^2\in[h^{2/Q_0},h^{2/Q_0}+h^2]$.
The window is available because $q<Q_0<3$.
For the large-coordinate first moments and dyadic class sum,
\[
 M(t)\le K\{h^{-7/30}+h^{-1/16}\},\qquad
 S(t)=\sum_j\sqrt{M_jV_j}\le K W,\qquad
 V(t)\le h^{7/16}.
\]
The head estimate uses $M(H>r)\le KW/\sqrt r$ and the gap
estimate uses $\int_t^r V(u)u^{-2}\,du$.

Here is the more precise class allocation needed at this step.
For a window of exponent $q$, the annular cost
$K(M^3/J^2+S^2/J+V^{3/2}/\sqrt J)$ is sufficient if
\begin{equation}\label{eq:weighted-annular-window}
 M\sigma^{q-4/3}L^{2/3}\to0,\quad
 S^2\sigma^{q-2}L\to0,\quad
 V^3L/\sigma^2\to0,\quad B\sigma^{2q-2}L\to0 .
\end{equation}
In fact put $U_0=\sigma^{2-2q}/L$,
$D_0=1+B+M^2L^{1/3}\sigma^{-2/3}+S^4L\sigma^{-2}$,
and $J=\lceil\sqrt{D_0U_0}\rceil$.
Then $D_0\ll J\ll U_0$. Applying
\eqref{eq:weighted-annular-cost}, the variance partition, and the
same retained prefix verifies the first two costs in
\eqref{eq:weighted-window-costs}. The third follows by multiplying
$w\sqrt{JL}/\sigma$ by
$\sigma J^{1/6}/\sqrt L\le\sigma^{(4-q)/3}L^{-2/3}$.
This proves the criterion.

For the displayed crossing parameters,
$(q-4/3)/Q_0=31/90>7/30,1/16$,
$(q-2)/Q_0=1/30$, and $3a/2-2/Q_0=91/240$.
Also $B=O(\log(1/t))=O(L)$. All conditions in
\eqref{eq:weighted-annular-window} hold with strict power margins,
giving the second contradiction.
\end{proof}

\section{A quantitative probability variation}\label{sec:weighted-packets}

Write the nonnegative full-law contact defect as
\[
 d(u)=-Q(u)=F(z)-\phi(z)u-Au^2+C|u|^3-F(z-u).
\]
To control inverse convolution without a continuity hypothesis, fix
$m>3$ and use
\begin{equation}\label{eq:weighted-envelope}
 E_m(h)=h^3+\sup_{r\ge h}(h/r)^m\sup_{|u|\le r}d(u).
\end{equation}
Then $h^3\le E_m(h)\le K$ for $h\le1$ and
$E_m(th)\le t^mE_m(h)$ for $t\ge1$. The added $h^3$ merely
includes the cubic polynomial in the same bound.

\begin{lemma}[Actual same-sign contacts]\label{lem:weighted-pair}
Translate a modal raw state of each of two distinct original
coordinates to zero, writing its law as
$K_i=(1-\theta_i)\delta_0+\eta_i$, $\theta_i\le1/4$.
Let the diameters be at most $Lh$, and choose actual states
$x=\varepsilon u$, $y=\varepsilon v$, with
$\kappa h\le u,v\le Lh$. Under
Lemma~\ref{lem:weighted-scalar}, uniformly as $h\to0$,
\[
 d(x)\le K\theta_1E_m(h),\quad
 d(y)\le K\theta_2E_m(h),\quad
 d(x+y)\le K(\theta_1+\theta_2)E_m(h).
\]
The bounds do not involve the probabilities of the selected states.
\end{lemma}
\begin{proof}
Write $T_a f(u)=f(u+a)$ and
$L_i=(T_{x_i}-I)K_i^{-1}$. The inverse is the absolutely
convergent series
$(1-\theta_i)^{-1}\sum_{n\ge0}[-T_{\eta_i}/(1-\theta_i)]^n$.
A term translated by $n$ raw states costs at most
$K(n+1)^mE_m(h)$ against $d$. Summing the geometric coefficients
proves that removing one inverse changes a finite difference by
$O(\theta_iE_m(h))$. Against $g(u)=|u|^3$ the same bound has
$h^3$ in place of $E_m(h)$.

Let $m_i$ be the raw mean and let $r_i,s_i$ be the first and
second derivatives of its centered cubic moment when mass moves
from zero to $x_i$. Direct differentiation gives
\[
 \begin{split}
 v_i'&=x_i^2-2m_ix_i,\\
 r_i&=|x_i-m_i|^3-|m_i|^3
             -3x_i\E[(Y_i-m_i)|Y_i-m_i|],\\
 s_i&=-6x_i\{(x_i-m_i)|x_i-m_i|+m_i|m_i|\}
                              +6x_i^2\E|Y_i-m_i|.
 \end{split}
\]
The exact first score is
$L_if(0)=-\phi x_i-Av_i'+Cr_i$, where $f(u)=F(z-u)$.
Since $L_i u=x_i$, $L_i u^2=v_i'$, this is
$L_id(0)=C(L_ig(0)-r_i)=O(\theta_i h^3)$.
Removing the inverse proves the individual bounds.

Keep the raw event fixed by compensating the centered threshold
by minus the changing mean. The original probability Hessian is
negative semidefinite, however small the two feasible probability
intervals. Divide its $ij$ entry by $x_ix_j$, a congruence.
With $a_i=v_i'/x_i$, $b_i=r_i/x_i$, its exact entries are
\[
 \begin{split}
 M_{ii}&=-3D+\phi''a_i-L_0a_i^2+3Ca_ib_i-Cs_i/x_i^2,\\
 M_{12}&=-3D+\tfrac{\phi''}2(a_1+a_2)-L_0a_1a_2
       +\tfrac{3C}2(a_1b_2+a_2b_1)
       +\{CL_1L_2g(0)-L_1L_2d(0)\}/(xy).
 \end{split}
\]
These follow by substitution in \eqref{eq:prob-hessian}; only
probabilities and normal normalization are differentiated.
Uniformly,
$a_i=x_i+O(\theta_i h)$,
$b_i=\varepsilon|x_i|^2+O(\theta_i h^2)$, and
$s_i/x_i^2=-6|x_i|+O(\theta_i h)$.
The double cubic difference is
$L_1L_2g(0)=3uv(u+v)+O((\theta_1+\theta_2)uvh)$:
its derivative under translation is bounded by $6|xy|$.
The double defect difference is
$d(x+y)-d(x)-d(y)+O((\theta_1+\theta_2)E_m(h))$.

Consequently the deterministic part is the symmetric kernel
\[
 B_\varepsilon(u,v)=-3D+c_\varepsilon(u+v)-L_0uv
                                      +\tfrac{3C}2uv(u+v).
\]
The diagonal errors are $O(\theta_i h)$ and the off-diagonal
upper error is $K(\theta_1+\theta_2)E_m(h)/(uv)$; the omitted
term $-d(x+y)/(uv)$ is nonpositive. Test on
$(1+t,-1)$, $t=c_\varepsilon(u-v)/(3D)$. Its two entries have
opposite signs, so that upper off-diagonal bound gives a lower
bound on the quadratic form. Direct expansion gives
\[
 (1+t,-1)B_\varepsilon(1+t,-1)^{\mathsf T}
  =(u-v)^2\{c_\varepsilon^2/(3D)-L_0+O(h)\}.
\]
Every polynomial term retains $(u-v)^2$, including when the gaps
coincide. The strict scalar margin gives
$uv(u-v)^2\le K(\theta_1+\theta_2)E_m(h)$.
Now retain the defect and test on $(1,-1)$. The polynomial form is
$(u-v)^2[-L_0+3C(u+v)]$, and the defect contributes
$2d(x+y)/(uv)$. The preceding estimate absorbs the polynomial
term and proves the pair bound.
\end{proof}

\begin{lemma}[A selected-intensity bound]\label{lem:weighted-intensity}
Fix $\kappa,L>0$. Consider finitely many actual coordinates with
modal raw states zero, diameters at most $Lh$, and modal deficits
at most $\epsilon$. In one physical sign, select submeasures of
their actual nonmodal states in $[\kappa h,Lh]$, of total mass
$0<\lambda\le1$. For sufficiently small $h,\epsilon$,
\begin{equation}\label{eq:weighted-intensity}
 \epsilon E_m(h)\ge c\lambda h^4\min(\lambda,h),\qquad
 \lambda\le K\{\sqrt{\epsilon E_m(h)}\,h^{-2}
                              +\epsilon E_m(h)h^{-5}\}.
\end{equation}
The independent remainder and all unselected states are retained.
\end{lemma}
\begin{proof}
Denote the selected submeasures by $\nu_i$, their masses by $r_i$,
and their normalized positive mark law by
$\mu=\lambda^{-1}\sum\nu_i$. Write
$a=\int x\,d\mu$, $b=\int x^2\,d\mu$, $c_3=\int x^3\,d\mu$.
They have their actual orders $h,h^2,h^3$ and $a\ge\kappa h$.
For sign $\varepsilon$ put
$s=2c_\varepsilon a/(3D)=O(h)$, $\alpha=2-s$.
For $q>0$ small use the count distribution with generating function
\begin{equation}\label{eq:weighted-compound-pgf}
 B_q(t)=e^{\alpha q(t-1)}
                         [1+sq(t-1)+q(t-1)^2].
\end{equation}
This is a probability law. To check positivity, its bracket is
$1+(1-s)q-\alpha qt+qt^2$. With $x=\alpha q$, its zeroth and
first coefficients are positive, the latter being
$e^{-x}\alpha q^2(1-s)$. For $n\ge2$, after extracting
$e^{-x}x^{n-2}q/(n-2)!$, the coefficient is
$1-\alpha^2q/n+\alpha^2q^2(1-s)/(n(n-1))>0$.
Here $0\le s\le1$ and a fixed sufficiently small upper bound on
$q$ suffices.

Append one centered coordinate
$\varepsilon(\sum_{j=1}^{N_q}Z_j-2qa)$, with iid marks of law
$\mu$. In original coordinate $i$, withdraw fraction
$\alpha q/\lambda$ of $\nu_i$ to its own modal state.
All laws stay positive for
$0\le q\le q_0\lambda$, with a fixed $q_0>0$.
Take the appended coordinate independently of the modified coordinates.

We bound the full convolution error, since an uncontrolled cubic
Taylor remainder would destroy the estimate. In the weighted
convolution norm
$\|\xi\|_m=\int(1+|u|/h)^m\,|\xi|(du)$, put
$D_\mu=T_{\varepsilon\mu}-I$ and
$A_i=\lambda^{-1}(T_{\varepsilon\nu_i}-r_iI)K_i^{-1}$.
The modal series gives
\[
 \sum_i\|A_i\|_m\le K,\quad
 \|D_\mu-\sum_iA_i\|_m\le K\epsilon,\quad
                 \sum_i\|A_i\|_m^2\le K\epsilon/\lambda.
\]
The relative convolution operator is
\[
 {\cal P}(q)={\cal R}(q)[I+q(sD_\mu+D_\mu^2)],\quad
 {\cal R}(q)=\prod_i(I-\alpha qA_i)e^{\alpha qD_\mu}.
\]
All factors commute. Its logarithmic derivative is
$\alpha(D_\mu-\sum A_i)
-\alpha^2q\sum A_i^2(I-\alpha qA_i)^{-1}$.
Thus $\|{\cal R}'(0)\|_m\le K\epsilon$ and
$\sup_{q\le q_0\lambda}\|{\cal R}''(q)\|_m\le K\epsilon/\lambda$,
whence
\begin{equation}\label{eq:weighted-relative-error}
 \|{\cal P}(q)-I-q{\cal P}'(0)\|_m
                                      \le K\epsilon q^2/\lambda.
\end{equation}
Against $d$ this costs $K\epsilon q^2E_m(h)/\lambda$.

The moments of the fresh law are also exact. Its first three
factorial count moments are
$2q$, $2q+(4-s^2)q^2$, and
$6\alpha q^2+(8-6s^2+2s^3)q^3$.
Let $P_0=\Pp(N_q=0)$. Since every positive marked value is at
least $\kappa h>2qa$, expansion on the zero and positive events gives
\[
 \begin{split}
 \Var(\sum Z_j)&=2q(b+a^2)-s^2a^2q^2,\\
 \E|\sum Z_j-2qa|^3
 &=q(2c_3+6ab)-q^2(3s^2ab+6sa^3)
                           +q^3(2s^3+16P_0)a^3 .
 \end{split}
\]
Put $e_1=sa$, $v_1=sb+2a^2$, $r_1=sc_3+6ab$.
The net raw mean change is exactly $\varepsilon e_1q$.
For the total variance and additive third-moment quantity,
\[
 \begin{split}
 V(q)-1&=q(v_1+O(\epsilon h^2))
                 -q^2(e_1^2+O(\epsilon h^2/\lambda)),\\
 {\cal B}'(0)&=r_1+O(\epsilon h^3),\\
 {\cal B}(q)-\beta-q{\cal B}'(0)
 &=-3e_1v_1q^2+(2s^3+16P_0)a^3q^3
                               +O(\epsilon q^2h^3/\lambda).
 \end{split}
\]
For the old coordinates these errors follow by differentiating
their affine means: their summed second cubic derivatives are at
most $K\lambda^{-2}\sum r_i^2h^3\le K\epsilon h^3/\lambda$.
The term $16P_0a^3q^3$ is retained as the cost of centering the
new coordinate.

Evaluate at centered threshold $z-\varepsilon e_1q$ so that the
raw event stays fixed. Every old first score is exactly zero.
The fresh first derivative is the insertion of two marks, hence
$J'(0)=-P_2$, where
\[
 P_2=\iint d(\varepsilon(x+y))\,d\mu(x)d\mu(y)
                    \le K\epsilon E_m(h)/\lambda.
\]
Lemma~\ref{lem:weighted-pair} bounds distinct-coordinate
selections; selections from the same original coordinate have
weight at most $\sum r_i^2/\lambda^2\le\epsilon/\lambda$ and
use the envelope. Thus a small fresh gradient is charged, not
set to zero.

For completeness, the smooth normalization after removing the
polynomial part of the CDF is
\[
 {\cal N}(e,v)=(1+v)^{3/2}
  [F(z)-\phi e-A(v+e^2)-\Phi((z-e)/\sqrt{1+v})]-D .
 \]
Its gradient at zero vanishes and its Hessian has entries
$-3D,\phi''(z)/2,-L_0$. Combining the preceding exact moments,
\eqref{eq:weighted-relative-error}, and the factor $(1+v)^{3/2}$
on the remaining cubic term gives
\[
 \begin{split}
 J(q)\ge -qP_2+\frac{q^2}{2}
   \{-3De_1^2+2c_\varepsilon e_1v_1-L_0v_1^2+3Cv_1r_1\}\\
               -K\{\epsilon q^2E_m(h)/\lambda+q^3h^3\}.
 \end{split}
\]
The product term involving $v_1P_2$ is absorbed in the displayed
error. The quadratic third-moment term supplies $6Ce_1v_1$ in the
Hessian; together with the normal mixed term this is
$2c_\varepsilon e_1v_1$. All remaining smooth terms are covered
by $q^3h^3$ or $\epsilon q^2E_m(h)/\lambda$, since $E_m(h)\ge h^3$.
Now $e_1=c_\varepsilon(2a^2)/(3D)$,
$v_1=2a^2+O(h^3)$, and $r_1=O(h^3)$. The braces are
$4a^4(c_\varepsilon^2/(3D)-L_0)+O(h^5)\ge c h^4$.
Choose $q=c_0\min(\lambda,h)$ with a small fixed $c_0$.
The centering cost is absorbed and $J(q)\le0$ proves
\eqref{eq:weighted-intensity}.
\end{proof}

\section{From grouping rates to weighted summability}
\label{sec:weighted-iteration}

\begin{lemma}[A grouping rate implies weighted variance summability]
\label{lem:weighted-rate-transfer}
Suppose that after deletion of any selected set of original coordinates,
the remaining coordinates can, for every large $J$, be partitioned
into at most $KJ$ groups with cost $w\le KJ^{-1-\theta}$,
where $\theta>0$ is fixed and the constant is independent of the
selected set. Then
\[
 V(h)=o(h^p)\quad\left(1<p<\frac{1+2\theta}{1+\theta}\right),
 \qquad
 W_s<\infty\quad\left(0<s<\frac{1+2\theta}{1+\theta}\right).
\]
\end{lemma}
\begin{proof}
If $V(h)\ge c h^p$ along a sequence, select a finite original
block with radii at most $h$ and variance $\sigma^2\asymp h^p$,
by stopping when a fixed smaller multiple of $h^p$ is reached.
Its variance overshoot is at most $h^2=o(h^p)$. Its independent
remainder is denoted by $A_0$ and has CDF $G$.
The majorant in Lemma~\ref{lem:weighted-majorant} has bounded
coefficients and $b+\phi(z)=o(1)$, while
Lemma~\ref{lem:weighted-normal} gives
$\E Q_U(\sigma Z)=O(\sigma^3+h)=O(h)$.

Choose fixed exponents
\[
 \frac{p(1+\theta)}{1+2\theta}<\lambda<1,\qquad
        \frac{\lambda}{1+\theta}<j<2\lambda-p,
 \qquad \ell=h^\lambda,\quad J=\lceil h^{-j}\rceil .
\]
All intervals are nonempty. In particular $\ell=o(\sigma)$,
$N\ell^2\log(1/h)=o(\sigma^2)$, and $w,h,\sigma^3=o(\ell)$.
Group the unselected remainder, append one ordinary
$N(0,\sigma^2)$ coordinate, and phase-round these $N+1$ coordinates
with total phase $z$ as in \eqref{eq:weighted-phase-identity}.
If $W_N$ is the sum of $N$ independent triangular noises, that
identity gives exactly
\[
 \E_{\rm phase}F_{\rm rounded}(z)
     =\E_{0\le u\le\ell}F_{A_0+\sigma Z+W_N}(z+u).
\]
The right-hand average is uniform in $u$.

We check the positive slack under this additional noise. Hoeffding's
elementary exponential-moment bound gives
$\E e^{tW_N}\le e^{N\ell^2t^2/2}$. If $g_u$ is the density of
$\sigma Z+W_N-u$, then
\[
 \frac{g_u(\sigma x)}{\phi_\sigma(\sigma x)}
 \le\exp\{\ell|x|/\sigma+N\ell^2x^2/(2\sigma^2)\}\le2
\]
on $|x|\le\sqrt{160\log(1/h)}$, uniformly for $0\le u\le\ell$.
The tails cost $o(h)$ against $Q_U\le K(1+|u|^3)$ by the same
subgaussian bound and sixth moments. Hence
$\sup_{0\le u\le\ell}\E Q_U(\sigma Z+W_N-u)=O(h)$.
For $R_u=\sigma Z+W_N-u$, the exact polynomial identity is
\[
 \begin{split}
 F_{A_0+\sigma Z+W_N}(z+u)-F(z)
 =-bu+c(N\ell^2/6+u^2)
 +C\{\E|R_u|^3-\E|U|^3\}
 +\E Q_U(U)-\E Q_U(R_u).
 \end{split}
\]
Its last positive term may be retained as zero in a lower bound.
Every term apart from $-bu$ is bounded below by
$-K(h+\sigma^3+N\ell^2+\ell^2)=o(\ell)$.
The phase-averaged CDF gain is thus
$\phi(z)\ell/2-o(\ell)$.

For every phase the variance lies between $1$ and
$1+(N+1)\ell^2/4=1+o(\ell)$. Its third-moment sum is at most
\[
 \beta+w+K\sigma^3+
        K\ell^2(\sqrt N+\sigma)+K(N+1)\ell^3=\beta+o(\ell).
\]
The removed third-moment sum was nonnegative; the upper bound
charges the new Gaussian and every grouped and rounded coordinate.
Normalization costs $o(\ell)$ uniformly. Some phase therefore has
positive objective, a contradiction. Finally
\eqref{eq:weighted-Tonelli}, with any larger available exponent
$p>s$, proves the assertion about $W_s$.
\end{proof}

\begin{lemma}[An initial polynomial grouping rate]
\label{lem:weighted-initial-rate}
Under the hypotheses of Lemma~\ref{lem:weighted-shape}, the uniform
grouping rate in Lemma~\ref{lem:weighted-rate-transfer} holds with
$\theta=1/80$. Consequently $W_s<\infty$ for every $s<82/81$ and
$\sum_i\E|X_i|<\infty$.
\end{lemma}
\begin{proof}
We have $W_{3/4}<\infty$ by Lemma~\ref{lem:weighted-subunit}.
Among coordinates with $v_i\le\delta$ and $H_i>r_\delta=\delta^{1/20}$,
balanced extremes force a triple's middle state to have modal
deficit $\theta_i\le K\delta/H_i^2$. Indeed the extreme contributions
to variance bound their total probability, and the centered middle
state has magnitude at most $\theta_iH_i$. All its nonmodal raw
gaps are therefore comparable to $H_i$. For two-point laws,
$v_i=\theta_i(1-\theta_i)H_i^2$ gives the same conclusion.
The finitely many exceptional coordinates are absent for small
$\delta$.

In each dyadic shell $h<H_i\le2h$, apply
Lemma~\ref{lem:weighted-intensity} separately to the two physical
signs, with $\epsilon\le K\delta/h^2$ and $E_m(h)\le K$.
It gives
$\lambda\le K(\sqrt\delta\,h^{-3}+\delta h^{-7})$.
For $h\ge r_\delta/2$ this tends to zero uniformly. If a sign's
total intensity reached one, select fractions totaling one to
obtain a contradiction. The entire sign intensity is therefore
less than one and obeys the same bound. Since $v_i\le\theta_iH_i^2$,
summing shells gives
\[
 {\cal W}(\delta):=\sum_{v_i\le\delta}v_i/\sqrt{H_i}
 \le W_{3/4}r_\delta^{1/4}
       +K(\sqrt\delta\,r_\delta^{-3/2}
                                  +\delta r_\delta^{-11/2})
 \le K\delta^{1/80}.
\]
Now remove any selected set. Retain all remaining $v_i>J^{-1/2}$
individually, at most $\sqrt J$ of them. At the fine cutoff
$r=J^{-2/3}$, the bounds from the full representation are
$M(r)\le KW_{3/4}r^{-1/4}$ and $V(r)\le W_{3/4}r^{3/4}$.
The remaining annular class sum is at most $K{\cal W}(J^{-1/2})$.
Annular and variance grouping give $O(J)$ groups and
\[
 Jw\le K\{M(r)^3/J+{\cal W}(J^{-1/2})^2
                                     +V(r)^{3/2}\sqrt J\}
       \le K(J^{-1/2}+J^{-1/80}+J^{-1/4}).
\]
These bounds also hold after deleting coordinates, so the constant is
uniform over selected sets. Lemma~\ref{lem:weighted-rate-transfer} proves the
weighted assertion, and \eqref{eq:weighted-shape-bound} with any
$s>1$ gives first-moment summability.
\end{proof}

\begin{lemma}[The envelope on a dyadic diameter interval]
\label{lem:weighted-rich-shell}
Suppose $W_s<\infty$ for some $s>1$. If a selected block
from $h<H_i\le2h$ has variance
$\sigma^2\ge c h^{2-\eta}$, where $\eta>0$ is fixed, then
$E_m(h)\le Kh$ for every fixed $m>\max(3,2/\eta)$.
\end{lemma}
\begin{proof}
Its coordinate radii are at most $2h$,
$h/\sigma=O(h^{\eta/2})$, and
$\sigma^2\le(2h)^sW_s=o(h)$. The majorant has
$b+\phi(z),c+A=O(\sigma)$ and
$\E Q_U(\sigma Z)=O(h)$.
Set $L=\sqrt{160\log(1/h)}$, $r_*=\sigma/L$.
For $|t|\le r_*$, the Gaussian density ratio on
$|u|\le\sigma L$ is bounded by $e^{1+1/(2L^2)}$, and its cubic
tails cost $o(h)$. The weighted Stieltjes estimate in
Lemma~\ref{lem:weighted-normal}, applied to $Q_U(\sigma x+t)$,
then gives
\[
 \sup_{|t|\le r_*}\E Q_U(U+t)\le Kh.
\]
Its error is bounded by
$K(h/\sigma)\{\sigma+(1+L^3)\E Q_U(\sigma Z+t)
 +(1+L^2)\phi(L)\}$, with both endpoints included; the actual
tail is $o(h)$ by Bernstein and sixth moments.
The exact averaged polynomial identity therefore gives
$0\le d(t)\le K(h+\sigma^2+\sigma^3)\le Kh$ for $|t|\le r_*$.
For larger radii use $d(t)\le K(1+|t|^3)$:
\[
 \sup_{r\ge r_*}(h/r)^m\sup_{|u|\le r}d(u)
                \le K(h/r_*)^m
                \le K h^{m\eta/2}L^m=o(h).
\]
Together with the added $h^3$ in \eqref{eq:weighted-envelope},
this proves the claim for the full envelope.
\end{proof}

\begin{lemma}[Coordinates with small variance or small $p_iH_i$]
\label{lem:weighted-sublevels}
If $1<s<12/7$ and $W_s<\infty$, then for $q=0,1/2$,
\begin{equation}\label{eq:weighted-variance-sublevels}
             \sum_{v_i\le\delta}v_iH_i^{-q}
                              \le K_s\delta^{(s-q)/6}.
\end{equation}
If all sufficiently late coordinates are two-point, write them
$X_i=\varepsilon_iH_i(B_{p_i}-p_i)$ with $p_i\le1/2$.
For $1<s<11/6$ with $W_s<\infty$,
\begin{equation}\label{eq:weighted-mean-sublevels}
             \sum_{p_iH_i\le\delta}p_iH_i
                              \le K_s\delta^{(s-1)/5}.
\end{equation}
\end{lemma}
\begin{proof}
For the first assertion choose
$0<\eta<2-7s/6$ and then $m>\max(3,2/\eta)$.
In the variance sublevel, each shell above
$r_0=L_0'\delta^{1/7}$ has modal deficit at most $K\delta/h^2$.
The global intensity bound from the preceding proof applies.
Choose the fixed constant $L_0'$ large enough that its right side
is below one at $h=r_0/2$. It bounds the full shell intensity
above that cutoff.

If a shell has total nonmodal intensity at least $h^{-\eta}$,
its actual variance is at least $ch^{2-\eta}$: the mode has mass
bounded below, the selected gaps are comparable to $h$, and
$\Var Y=\frac12\E(Y-Y')^2$ includes the mode--nonmode pairs.
Lemma~\ref{lem:weighted-rich-shell} gives $E_m(h)\le Kh$.
Selecting one unit of intensity in a sign and using
\eqref{eq:weighted-intensity} forces $\delta\ge c h^6$.
Thus above $r_1=L_1'\delta^{1/6}$ the shell intensity is less
than $h^{-\eta}$, with a sufficiently large fixed $L_1'$.
For small $\delta$, $r_1<r_0$. The low, middle and high ranges
have the following bounds:
\[
\begin{array}{c|c}
 H_i\le r_1 & K\delta^{(s-q)/6}\\
 r_1<H_i\le r_0 & K\delta^{(2-q-\eta)/7}\\
 H_i>r_0,\ v_i\le\delta&
       K\{\delta^{(5-2q)/14}+\delta^{(2-q)/7}\}.
\end{array}
\]
The low range uses $W_s$; the others sum geometric shell bounds
$h^{2-q}$ times their actual intensities. The chosen range of
$s,\eta$ makes the last two rows no larger than the first, proving
\eqref{eq:weighted-variance-sublevels}.

For a two-point coordinate put $m_i=p_iH_i$. On $m_i\le\delta$,
the modal deficit is $p_i\le\delta/h$ in a shell. The global
intensity bound is
$K(\sqrt\delta\,h^{-5/2}+\delta h^{-6})$, valid for the entire
shell above $r_0=L_0'\delta^{1/6}$ by the same unit-selection
argument. Choose $0<\eta<(11-6s)/5$. If the shell intensity is at least
$h^{-\eta}$, then
Lemma~\ref{lem:weighted-rich-shell} gives $E_m(h)\le Kh$, and selecting
one unit of intensity forces $\delta\ge c h^5$. Thus shell intensity is below $h^{-\eta}$
above $r_1=L_1'\delta^{1/5}$. The three contributions to
$\sum m_i$ are
\[
 K\delta^{(s-1)/5},\qquad
 K\delta^{(1-\eta)/6},\qquad
 K(\delta^{1/4}+\delta^{1/6}).
\]
For the first use $m_i=v_i/((1-p_i)H_i)\le2v_i/H_i$.
The stated range of $s,\eta$ makes all exponents at least
$(s-1)/5$. The finite exceptional prefix is absent from either
sublevel for sufficiently small $\delta$.
\end{proof}

\begin{proof}[Proof of Theorem~\ref{thm:weighted-variance}]
The elementary counting implication
\[
 \sum_{a_i\le\delta}a_i\le K\delta^b,\quad0<b<1
          \quad\Longrightarrow\quad
 \#\{i:a_i>\delta\}\le K'\delta^{b-1}
\]
follows by summing the dyadic intervals
$(2^k\delta,2^{k+1}\delta]$. For the general core apply it to
\eqref{eq:weighted-variance-sublevels} with $q=0$ and choose
$\delta=J^{-6/(6-s)}$. Retain all $v_i>\delta$, at most $KJ$
original coordinates. The remaining $W_{1/2}$ is at most
$KJ^{-(s-1/2)/(6-s)}$. Since the full first-moment sum is finite,
annular grouping above diameter $J^{-1}$ and variance grouping below
it give, after any selected deletion,
\[
 w\le K_s\{J^{-2}+J^{-1-2(s-1/2)/(6-s)}
                                  +J^{-1/2-3s/2}\}
       \le K_sJ^{-1-\theta_s},\qquad
                    \theta_s=\frac{2s-1}{6-s}.
\]
The class count is $O(\log J)$ and is included in $O(J)$.
Lemma~\ref{lem:weighted-rate-transfer} therefore sends every
available $s\in(1,12/7)$ to every
$\gamma<T(s)=(3s+4)/(s+5)$.
Initially an $s>1$ is available by
Lemma~\ref{lem:weighted-initial-rate}. Availability is downward
closed. If its supremum below $\sqrt5-1$ were $s_*$, choose
available $s$ close enough to $s_*$ that $T(s)>s_*$, a contradiction.
Here $T(s)>s$ exactly below the positive fixed point $\sqrt5-1$.
Every use fixes $s$ before the limit; no uniform constants or
conclusion at the endpoint are needed.

For the two-point tail apply the counting implication to
\eqref{eq:weighted-mean-sublevels} and choose
$\delta=J^{-5/(6-s)}$. The retained count is $O(J)$ and the
remaining endpoint mean is at most $KJ^{-(s-1)/(6-s)}$.
Group the coordinates according to the sign of their nonmodal displacement.
Within each sign class, apply the grouping estimate after reflection if
necessary, using the nonnegative raw variables $H_iB_{p_i}$. Their
negative first-moment term is zero.
Lemma~\ref{lem:weighted-grouping} gives
\[
 w\le K_sJ^{-2-3(s-1)/(6-s)}
       =K_sJ^{-1-\theta_s^{\rm B}},\qquad
                   \theta_s^{\rm B}=\frac{2s+3}{6-s}.
\]
The finite original prefix is retained individually. This is uniform
after any selected deletion. The new map is
$T_{\rm B}(s)=(3s+12)/(s+9)$, with positive fixed point
$\sqrt{21}-3<11/6$. Starting from any available $s>1$, the same
supremum argument proves the improved range. First-moment
summability already followed from
Lemma~\ref{lem:weighted-initial-rate}.
\end{proof}

\subsection{Regularity at the selected threshold}
\label{subsec:weighted-regularity}

\begin{proof}[Proof of Corollary~\ref{cor:weighted-regularity}]
The general grouping rate above has some $\theta>0$. If
$F(z+t)=F(z)+\phi(z)t+O(t^2)$ through both signs, take
$\ell=J^{-a}$ with $1<a<1+\theta$. The exact phase kernel
\eqref{eq:weighted-phase-identity} has average linear gain
$\phi(z)\ell/2$ and quadratic error $O(J\ell^2)=o(\ell)$.
Concentration pays its excursion outside the stated neighborhood.
Its actual grouping, rounding and normalization costs are $o(\ell)$,
as in the preceding phase proofs. This contradicts maximality.

More generally, suppose only the left expansion
$F(z)-F(z-h)=\phi(z)h+O(h^{1+\alpha})$ holds.
It is enough to take $0<\alpha\le1$, decreasing a larger exponent
if necessary. A grouping rate $w=O(J^{-1-\theta})$ excludes it
whenever $\alpha>1/(1+2\theta)$. Choose
\[
          \frac{1+\alpha}{2\alpha}<a<1+\theta,\qquad
          \ell=J^{-a}.
\]
Evaluate the rounded array at the lattice point $z-b_J\ell$,
where the integer $b_J=O(\sqrt{J\log J})$ is large enough that
$-b_J\ell-W_\ell\in[-2b_J\ell,0]$ except with probability
$O(J^{-L})$, for a fixed $L>a+1$.
Hoeffding's inequality supplies this event. The exact noise mean
and the left expansion give
\[
 \E_{\rm phase}F_{\rm rounded}(z-b_J\ell)
   =F(z)-\phi(z)b_J\ell+\phi(z)\ell/2+o(\ell),
\]
since $\ell^\alpha(J\log J)^{(1+\alpha)/2}\to0$.
The bounded CDF pays the exceptional event. Normal Taylor error
is $O(J\ell^2\log J)=o(\ell)$; the variance cost is
$O(J\ell^2)$, and the actual cubic cost is
$w+O(\sqrt J\ell^2+J\ell^3)=o(\ell)$.
Thus the same positive continuity correction remains.

For the two-point tail, the rate just proved has
$\theta=(2s+3)/(6-s)$ for every fixed available
$s<\sqrt{21}-3$. Its threshold is
$1/(1+2\theta)=(6-s)/(3s+12)$, whose limit is
$(\sqrt{21}-3)/6$. Choose $s$ sufficiently close to the endpoint.
\end{proof}
